\documentclass[11pt,a4paper]{article}
\usepackage[utf8]{inputenc}
\usepackage[T1]{fontenc}
\usepackage{amsmath,amsfonts,amssymb}
\usepackage{graphicx}
\usepackage{hyperref}
\usepackage{listings}
\usepackage{color}
\usepackage{booktabs}
\usepackage{float}
\usepackage{geometry}
\geometry{margin=1in}
\definecolor{zenblue}{RGB}{41,121,255}
\hypersetup{colorlinks=true,linkcolor=zenblue,urlcolor=zenblue,citecolor=zenblue}

\title{\textbf{Zen-Dub: A Pipeline for AI Voice Dubbing and Localization}\\
\large Technical Whitepaper v2025.03}
\author{Zach Kelling \\ Zen LM Research Team\\
\texttt{research@zenlm.org}\\
\href{https://papers.zenlm.org}{papers.zenlm.org}}
\date{March 2025}

\begin{document}
\maketitle

\begin{abstract}
Zen-Dub is a video-dubbing \emph{pipeline} assembled from existing, openly-licensed
components---not a single model trained from scratch, and not a ``Mixture of Distilled
Experts.'' It chains automatic speech recognition, machine translation, voice-cloning
text-to-speech, and audio-driven lip synchronization into a localization workflow that
preserves each speaker's voice identity and matches mouth motion. Voice cloning uses
Alibaba's \textbf{Qwen3-TTS}~\cite{qwen3tts} (Apache~2.0), with \textbf{CosyVoice~2}
~\cite{cosyvoice2} (Apache~2.0) as a streaming backup engine. Lip synchronization uses
\textbf{MuseTalk}~\cite{musetalk} (Lyra Lab, Tencent Music; MIT code), a real-time
latent-space inpainting model that regenerates the mouth region at 256$\times$256 and
30+\,fps. This whitepaper describes the pipeline's stages and the provenance and
license of each component. We do not report dubbing-quality benchmark numbers of our
own; capability claims are attributed to the upstream components rather than
re-measured here.
\end{abstract}

\tableofcontents
\newpage

%% -----------------------------------------------------------------------
\section{Introduction}
\label{sec:intro}

Professional voice dubbing is one of the most labor-intensive tasks in media localization.
For a 90-minute film, professional dubbing requires:
\begin{itemize}
\item Script translation and adaptation by localization experts
\item Casting native-language voice actors
\item Recording sessions in acoustic studios
\item Audio post-production (mixing, timing adjustment, EQ matching)
\end{itemize}

This pipeline costs \$50,000--\$500,000 per language for feature-length content and
takes 4--12 weeks per language. The result: most content is localized into fewer than
10 languages, leaving global audiences underserved.

Zen-Dub compresses this workflow into an automated multi-stage pipeline (a sequence of
existing models, not a single trained model), producing dubbed audio that preserves:

\begin{enumerate}
\item \textbf{Speaker voice identity} --- Timbre, pitch range, resonance, and speaking
      style characteristics extracted from source audio, cloned into the target language.
\item \textbf{Emotional cadence} --- Prosodic features (speech rate variation, emphasis
      patterns, pause placement) are detected in source audio and reproduced in target
      synthesis with language-appropriate phoneme mapping.
\item \textbf{Lip-sync timing} --- Phoneme duration in the target language synthesis
      is adjusted to match source mouth movement timing, achieving natural lip-sync
      without visible audio lag.
\item \textbf{Background audio preservation} --- Source non-speech audio (music,
      ambient sound, sound effects) is automatically separated and preserved in the
      output mix.
\end{enumerate}

%% -----------------------------------------------------------------------
\section{Architecture}
\label{sec:arch}

Zen-Dub is a multi-component pipeline. There is no single trained ``backbone'' and no
``Mixture of Distilled Experts''; each stage is a distinct, separately-licensed
component behind a stable interface, and the heavy lifting is done by upstream models
(Table~\ref{tab:components}). The pipeline consists of five stages:

\begin{enumerate}
\item Source audio processing (ASR + speaker characterization + source separation)
\item Machine translation (isochrony-aware)
\item Prosody transfer and lip-sync timing adaptation
\item Voice synthesis (cloning TTS: Qwen3-TTS or CosyVoice~2)
\item Lip synchronization (MuseTalk) and audio post-processing / mixing
\end{enumerate}

\begin{table}[H]
\centering
\caption{Pipeline components and their provenance. Zen-Dub integrates these; it does
not train them.}
\label{tab:components}
\begin{tabular}{lll}
\toprule
\textbf{Stage} & \textbf{Component} & \textbf{License} \\
\midrule
ASR + audio features & Whisper-family encoder & MIT \\
Translation & Open MT model (swappable) & per upstream \\
Voice clone (primary) & Qwen3-TTS~\cite{qwen3tts} & Apache 2.0 \\
Voice clone (backup) & CosyVoice~2~\cite{cosyvoice2} & Apache 2.0 \\
Source separation & Demucs & MIT \\
Lip-sync & MuseTalk~\cite{musetalk} & MIT (code) \\
\bottomrule
\end{tabular}
\end{table}

\subsection{Source Audio Processing}

\textbf{Automatic Speech Recognition (ASR)}: An off-the-shelf ASR model (Whisper-family)
transcribes source audio with word-level timestamps. We do not train this component;
recognition quality is the upstream model's.

\textbf{Speaker characterization}: A speaker-embedding network extracts a per-segment
speaker identity vector used to condition the cloning TTS. We use an existing speaker
encoder rather than a bespoke one.

\textbf{Background separation}: An existing source-separation model (Demucs, MIT)
isolates speech from music and ambient sound so the non-speech bed can be preserved in
the final mix.

\subsection{Isochrony-Aware Machine Translation}

Translation in dubbing differs from document translation: the output must be
semantically accurate, phonetically compatible with the source timing, and natural in
spoken form. Zen-Dub uses an off-the-shelf machine-translation model (swappable) and
applies an \emph{isochrony} constraint at decoding time so that the target length
matches the source. The constraint is a property of the pipeline, not of a custom
model:

\textbf{Isochrony constraints}: The translation length (in phonemes) must approximately
match the source length (in phonemes) to preserve lip-sync. A duration-aware
decoding procedure biases beam search toward translations with matching phoneme counts:

\begin{equation}
\text{score}(y | x) = \log p_{\text{MT}}(y | x) - \lambda \left| \frac{|y|_{\phi}}{|x|_{\phi}} - 1 \right|
\end{equation}

where $|y|_\phi$ and $|x|_\phi$ are the estimated phoneme counts of target and source
strings, and $\lambda = 2.0$ is the isochrony penalty.

\textbf{Spoken form optimization}: A post-edit module rewrites written translations
into spoken-language style: contractions, natural spoken discourse markers, and
abbreviation expansion appropriate to the target language's spoken register.

\textbf{Character voice preservation}: Named entities and proper nouns are handled
with a character name preservation dictionary, ensuring consistency of character
names across the dub.

\subsection{Prosody Transfer}

Prosody (speech rhythm, stress, intonation) is transferred from source to target
via a prosody encoder-decoder:

\begin{equation}
\hat{\mathbf{p}}_{\text{target}} = f_{\text{prosody}}\!\left(\mathbf{p}_{\text{source}}, \mathbf{e}_{\text{emotion}}, \mathbf{l}_{\text{target}}\right)
\end{equation}

where $\mathbf{p}_{\text{source}}$ is the source prosody representation (F0 contour,
energy envelope, duration sequence), $\mathbf{e}_{\text{emotion}}$ is the emotion
embedding, and $\mathbf{l}_{\text{target}}$ is a learned target language prosody prior.

The prosody decoder outputs word-level duration targets and phrase-level F0 contours
in the target language, respecting the phonological constraints of the target language
(e.g., tonal languages receive special treatment of F0 mapping).

\subsection{Lip-Sync Adaptation}

Lip-sync accuracy requires that the synthesized audio's viseme sequence align with
the source speaker's mouth movements at a coarse level ($\pm$80ms). The lip-sync
adapter modifies phoneme durations in the synthesis plan:

\begin{equation}
d_i^* = \arg\min_{d_i} \left\{ |d_i - \hat{d}_i| + \mu \sum_j | \tau_j^{\text{sync}} - \tau_j^{\text{source}} | \right\}
\end{equation}

where $d_i$ are phoneme durations, $\hat{d}_i$ are the natural prosody targets,
$\tau_j^{\text{sync}}$ are synthesized viseme boundary times, and $\tau_j^{\text{source}}$
are source mouth-open/close event times detected by a lightweight face motion detector.

\subsection{Voice Synthesis with Speaker Cloning}

Voice synthesis is delegated to an upstream zero-shot cloning TTS model---primarily
\textbf{Qwen3-TTS}~\cite{qwen3tts}, with \textbf{CosyVoice~2}~\cite{cosyvoice2} as a
streaming backup---conditioned on the extracted speaker reference, the translated
phoneme/text sequence, and the transferred prosody plan:

\begin{equation}
\hat{A} = \text{CloneTTS}\!\left(\text{text}_{\text{target}}, \mathbf{s}_{\text{speaker}}, \mathbf{p}_{\text{prosody}}\right)
\end{equation}

The cloning behavior (e.g.\ synthesis from a short reference) is a property of the
upstream model; per Alibaba's release, Qwen3-TTS clones from roughly three seconds of
reference audio. We cite this rather than re-measuring it. Zen-Dub does not train the
TTS model.

%% -----------------------------------------------------------------------
\section{Integration, Not Training}
\label{sec:training}

Zen-Dub performs \emph{no} end-to-end model training. There is no joint ASR+MT+TTS
multi-task objective, no 970,000-hour dubbing corpus, and no ``human evaluation loop''
driving fine-tuning---earlier drafts described such a training process, but it does not
exist and those claims have been removed. Each stage is an existing model used as
released; Zen-Dub's engineering is the glue: format conversion, the isochrony decoding
constraint, prosody/duration adaptation, the lip-sync timing alignment, and the final
remix under the preserved background bed.

Dataset and training details for the underlying models are documented by their
respective upstream sources (Qwen3-TTS~\cite{qwen3tts}, CosyVoice~2~\cite{cosyvoice2},
MuseTalk~\cite{musetalk}).

%% -----------------------------------------------------------------------
\section{Evaluation}
\label{sec:eval}

We do not present dubbing-quality benchmark tables (speaker-similarity MOS, BLEU,
lip-sync accuracy, throughput) in this whitepaper. Earlier versions contained such
tables with specific numbers (e.g.\ ``MOS 4.3'', ``BLEU 42.3'', ``87.4\% lip-sync''),
but those numbers were not the result of a controlled evaluation of this pipeline and
have been removed rather than presented as fact.

For component-level quality, readers should consult the upstream reports:
\begin{itemize}
\item \textbf{Voice cloning}: Qwen3-TTS~\cite{qwen3tts} and CosyVoice~2~\cite{cosyvoice2}.
\item \textbf{Lip synchronization}: MuseTalk~\cite{musetalk}, which reports real-time
      256$\times$256 generation at 30+\,fps and competitive visual fidelity / lip-sync
      accuracy.
\end{itemize}

A sibling report, \emph{Zen Live-Dub} (newsroom), contains directly-measured,
independently-verified numbers for a closely-related license-clean dubbing pipeline
(lip-sync throughput, cross-lingual clone similarity, watermark survival); we point
readers there for measured figures instead of repeating unverified ones here.

%% -----------------------------------------------------------------------
\section{Related Work}
\label{sec:related}

\subsection{Neural Dubbing}

Prior work on automatic dubbing focused primarily on isochrony constraints in
machine translation~\cite{zetterholm2004} and prosody transfer~\cite{Federico2020}.
Neural dubbing systems combine ASR, MT, voice-cloning TTS, and speaker adaptation.
Zen-Dub follows this pipeline pattern, assembling existing models and adding
isochrony-aware decoding, prosody/timing adaptation, and MuseTalk-based lip
synchronization, rather than training a new end-to-end model.

\subsection{Voice Cloning}

Zero-shot voice cloning from short reference audio has advanced rapidly. Zen-Dub's
synthesis stage is the upstream Qwen3-TTS model~\cite{qwen3tts} (with CosyVoice~2
~\cite{cosyvoice2} as a backup), both Apache-2.0; we do not introduce a new vocoder.

\subsection{Prosody Transfer}

Cross-lingual prosody transfer is complicated by language-specific intonation systems.
Work on prosody normalization and language-aware F0 transfer~\cite{prosody} has shown
that speaker-independent prosody components (emphasis, rhythm) transfer better than
language-dependent ones (tonal patterns in tonal languages).

%% -----------------------------------------------------------------------
\section{Conclusion}
\label{sec:conclusion}

Zen-Dub is a localization \emph{pipeline} that chains existing, openly-licensed
models---Qwen3-TTS (with CosyVoice~2 as backup) for voice cloning and MuseTalk for lip
synchronization---behind ASR, isochrony-aware translation, and prosody/timing
adaptation. Its contribution is integration and license-clean assembly, not a novel
trained model and not the ``Zen MoDE'' architecture earlier drafts claimed.

The practical value is that a multilingual dubbing workflow can be stood up from
permissively-licensed parts, with each component's quality attributable to its
upstream source. We deliberately make no quantitative quality claim of our own in this
document; measured figures for a closely-related pipeline appear in the Zen Live-Dub
report.

Future work targets emotional consistency across scene cuts, multi-speaker handling for
crowd scenes, and singing-voice dubbing for musical content.

\section*{Acknowledgments}

We thank the upstream teams whose openly-licensed models this pipeline integrates: the
Qwen team at Alibaba Cloud (Qwen3-TTS), FunAudioLLM (CosyVoice~2), and Lyra Lab at
Tencent Music (MuseTalk).

\bibliographystyle{plain}
\begin{thebibliography}{9}

\bibitem{qwen3tts}
Qwen Team, Alibaba Cloud,
``Qwen3-TTS: Open-Source Streaming Text-to-Speech with Voice Cloning,''
\url{https://github.com/QwenLM/Qwen3-TTS}, 2026. Apache 2.0.

\bibitem{cosyvoice2}
Z. Du et al. (FunAudioLLM, Alibaba),
``CosyVoice 2: Scalable Streaming Speech Synthesis with Large Language Models,''
\textit{arXiv:2412.10117}, 2024. Apache 2.0.

\bibitem{musetalk}
Y. Zhang et al. (Lyra Lab, Tencent Music Entertainment),
``MuseTalk: Real-Time High-Fidelity Video Dubbing via Spatio-Temporal Sampling,''
\textit{arXiv:2410.10122}, 2024. Code: MIT.

\bibitem{zetterholm2004}
E. Zetterholm,
``Voice Imitation: A Phonetic Study of Perceptual Illusions and Acoustic Success,''
\textit{Lund University}, 2004.

\bibitem{Federico2020}
M. Federico et al.,
``From Speech to Speech Translation,''
\textit{ICASSP}, 2020.

\bibitem{vits2}
J. Kong et al.,
``VITS2: Improving Quality and Efficiency of Single-Stage Text-to-Speech with Adversarial Learning and Architecture Design,''
\textit{Interspeech}, 2023.

\bibitem{prosody}
C. Wan et al.,
``Prosody Transfer in Neural Machine Translation for Dubbing,''
\textit{ACL}, 2023.

\end{thebibliography}

\end{document}
